% paper/sections/15_conclusion.tex
\section{まとめと今後の課題}
\label{sec:conclusion}

% ═══════════════════════════════════════════════════════════════════════════════
\subsection{本稿の貢献のまとめ}
\label{sec:contribution_summary}
% ═══════════════════════════════════════════════════════════════════════════════

本稿では，気液二相流の高精度数値シミュレーションに必要な数学的・数値的基盤を，
学部生・大学院生を主な読者として体系的に解説した．
本稿の技術的中核は\textbf{寄生流れの根本的抑制}にあり，
（i）CCD $\Ord{h^6}$ を起点とする高精度微分演算子群（CCD・DCCD・UCCD6・FCCD・面ジェット）による曲率と勾配の同時高精度化，
（ii）\textbf{分相 PPE}＋HFE による界面延長＋欠陥補正（DC $k=3$）の組合せによる変密度 PPE 収束性の回復，
（iii）Balanced-Force 七原則と FCCD 面ジェット sub-system による
コロケート格子のチェッカーボード抑制，
の \textbf{三本柱}でこれを実現する（§\ref{sec:intro}，§\ref{sec:roadmap}，§\ref{sec:strategy}）．
各章の設計方針を整理する前に，「なぜこの組合せが必要か」を再確認する．

\begin{center}
\renewcommand{\arraystretch}{\PaperTableArrayStretch}
\begin{tabular}{>{\raggedright}p{58mm}>{\raggedright\arraybackslash}p{82mm}}
\hline
\textbf{数値的困難} & \textbf{本稿の対処手法} \\
\hline
界面の質量損失       & CLS 法 + 6 段階再初期化（第\ref{sec:levelset}章，§\ref{sec:reinit}，§\ref{sec:cls_stages}）\\
低次移流による界面にじみ & DCCD + UCCD6 + TVD-RK3（§\ref{sec:dccd_role}，§\ref{sec:uccd6_def}，§\ref{sec:time_int}）\\
曲率計算誤差 $\to$ 寄生流れ & CCD $\Ord{h^6}$ + HFE Hermite 場延長（第\ref{sec:CCD}章，§\ref{sec:hfe}）\\
チェッカーボード圧力振動 & Balanced-Force 七原則 + FCCD 面ジェット sub-system（§\ref{sec:balanced_force}，§\ref{sec:fccd_bf_sub_system}）\\
変密度 PPE の収束悪化 & 分相 PPE + HFE + 欠陥補正の収束判定（§\ref{sec:split_ppe_recovery}，§\ref{sec:hfe}，§\ref{sec:defect_correction_ppe}）\\
距離関数品質の劣化と位相変化 & Ridge--Eikonal 補助場再初期化（§\ref{sec:ridge_eikonal}，§\ref{sec:ridge_aux_field}）\\
格子粗さ・壁接触による界面解像不足 & 界面適合非一様格子 + Ridge--Eikonal 非一様補正 D1--D4 + 壁閉包 W（第\ref{sec:grid}章，§\ref{sec:ridge_eikonal_nu_grid}）\\
高密度比下の圧力ジャンプ整合 & pressure-jump 型 FCCD-PPE + 毛管 cochain 値域射影（§\ref{sec:split_ppe_pressure_jump_formulation}，§\ref{sec:split_ppe_capillary_range_projection}）\\
アクティブ幾何毛管分解の採用境界 & 幾何セル分率 $q_C$，互換射影，圧力反力部分空間 $R_p$，Hodge 逆適用済み面補間，移動格子面共鎖，圧力履歴座標，近似誤差境界（§\ref{sec:ao_fast_state_space}）\\
起動時の格子↔界面循環依存 & DCCD ブートストラップ 6 段階（§\ref{sec:bootstrap}）\\
\hline
\end{tabular}
\end{center}

% ─────────────────────────────────────────────────────────────────────────────
\subsubsection{三本柱の設計方針と各章の技術的貢献}
% ─────────────────────────────────────────────────────────────────────────────

以下，三本柱に沿って Part I（定式化）・Part II（数値）の各章の技術的貢献を整理する．
末尾に三本柱を結ぶ\textbf{統合層}（非一様格子拡張・統合アルゴリズム・ブートストラップ・Pure-FCCD DNS）をまとめる．

\paragraph{第一の柱：CCD $\Ord{h^6}$ を起点とする高精度微分演算子群と界面追跡}\hfill\null

\begin{description}
  \item[\textbf{第\ref{sec:governing}章：支配方程式系と One-Fluid 定式化・表面張力モデル}]
        2つの独立した Navier--Stokes 方程式系（液相・気相）を，
        ヘヴィサイド関数 $\psi$ を用いた一体化方程式に変換する
        One-Fluid 定式化を導出した（§\ref{sec:onefluid}）．
        界面ジャンプ条件 $[p]_\Gamma=\kappa/We$（Young--Laplace 則）の離散的取り扱いを
        CSF 体積力形式（付録~\ref{app:csf_full}）と HFE（§\ref{sec:hfe}）の
        両面から整理した．
        無次元化（Re，Fr，We の定義）と曲率公式の幾何学的導出（§\ref{sec:curvature}）により，
        後の章の数値実装の数学的基礎を確立した．

  \item[\textbf{第\ref{sec:levelset}章：Conservative Level Set 法と Ridge--Eikonal 再初期化}]
        標準 Level Set 法の質量損失問題を克服するため，
        $\psi = H_\varepsilon(-\phi) \in [0,1]$ を移流変数とする保存形の
        界面追跡法（CLS 法）を定式化した．
        $\psi \to \phi$ の解析的逆変換（§\ref{sec:newton_inversion}）と
        Ridge--Eikonal ハイブリッド再初期化（§\ref{sec:ridge_eikonal}，§\ref{sec:ridge_aux_field}）
        により，飽和帯と位相変化の両者で連続な距離関数復元を可能にした．
        単体検証 U4 では DGR の有効幅比が
        $\varepsilon_{\mathrm{eff}}/\varepsilon_*=2.0\to1.0000033$ に収束し，
        設計どおりの帯幅補正が実証された．

  \item[\textbf{第\ref{sec:CCD}章：CCD 法ファミリと面ジェット sub-system}]
        3点ステンシルで1階・2階微分を同時求解する CCD 法（Chu \& Fan 1998）を定式化し，
        係数導出根拠（付録~\ref{app:ccd_coef_derivation}）と
        楕円型境界値問題への仮想時間陰解法応用（§\ref{sec:ccd_elliptic}）を整理した．
        本稿はこれを基幹に\textbf{4 種の派生演算子}を導入する：
        散逸 CCD（DCCD; §\ref{sec:dccd_role}，スペクトルフィルタ
        $H(\xi)=1-4\varepsilon_d\sin^2(\xi/2)$ による $2\Delta x$ モード抑制），
        6次風上 CCD（UCCD6; §\ref{sec:uccd6_def}, §\ref{sec:uccd6_hyperviscosity}，
        運動量対流の節点中心ハイパー粘性），
        面中心 CCD（FCCD; §\ref{sec:fccd_def}，$\Ord{h^4}$ 設計），
        面ジェット API（§\ref{sec:face_jet_def}, §\ref{sec:face_jet_roles}，
        $J_f(u)=(u_f,u'_f,u''_f)$ を 4 経路（圧力勾配・表面張力・HFE 上流点・GFM ジャンプ）で共有）．
        単体検証 U1 で CCD $\Ord{h^{6.0}}$，DCCD $H(\pi)=0$，FCCD $\Ord{h^{4.00}/h^{4.00}}$，
        UCCD6 $\Ord{h^{7.00}}$ のスーパー収束を確認した．

  \item[\textbf{第\ref{sec:reinit}章：CLS 再初期化パイプラインと 6 段階手順}]
        Ridge--Eikonal を Stage D に組み込んだ
        CLS の 6 段階離散サイクル（A: FCCD 移流，B: $\psi$ クランプ，
        C: ロジット逆変換，D: Godunov Eikonal 再初期化，
        E: DGR 帯幅補正，F: 曲率更新）を §\ref{sec:cls_stages} で展開し，
        質量保存・幾何精度・質量補正の責務分離を実現した．

  \item[\textbf{第\ref{sec:scheme_per_variable}章：変数別スキームと 3 層粘性離散化}]
        CLS $\psi$ には FCCD 保存型 flux 発散（DCCD $\varepsilon_d=0.05$ 併用），
        運動量 $(u,v)$ には UCCD6，
        粘性項には 3 層 CCD（Layer A/B/C）＋ Crank--Nicolson 半陰解法を割り当て，
        $\rho/\mu$ の補間規則（$\psi$ 線形）を明文化した（§\ref{sec:scheme_per_variable}）．

  \item[\textbf{第\ref{sec:time_int}章：時間積分（TVD-RK3 / AB2+IPC / CN ADI）}]
        CLS と移流に TVD-RK3，NS には AB2+IPC 投影による
        \textbf{全体時間精度 $\Ord{\Delta t^2}$} を採用し，
        粘性 ADI Crank--Nicolson と毛管波制約 $\Delta t_\sigma \propto \Delta x^{3/2}$ を導出した．
        単体検証 U8 で TVD-RK3 $\Ord{\Delta t^{3.01}}$，AB2 $\Ord{\Delta t^{2.04}}$，
        CN $\Ord{\Delta t^{2.00}}$ を実測した．
\end{description}

\paragraph{第二の柱：分相 PPE + HFE + 欠陥補正による変密度 PPE 収束性の回復}\hfill\null

\begin{description}
  \item[\textbf{第\ref{sec:ppe_solve}章：CCD-Poisson と 分相 PPE / HFE / 欠陥補正}]
        変密度 IPC Projection 法から
        変係数 PPE $\nabla\cdot(\rho^{-1}\nabla p)=(1/\Delta t)\nabla\cdot\bu^*$ を導出した．
        本稿の主たる設計戦略は\textbf{分相 PPE}（§\ref{sec:split_ppe_recovery},
        §\ref{sec:varrho_ppe_limitation}）であり，
        各相内で定数密度 Poisson を独立に解き，
        界面ジャンプ条件 $j_{gl}=p_g-p_l=-\sigma\kappa_{lg}$ を
        \textbf{Hermite 場延長 HFE}（§\ref{sec:hfe}）が
        Aslam Extension PDE の定常解として
        Hermite 三つ組 $(f,f',f'')$ から仮想時間反復なしに直接構成する．
        さらに各相に\textbf{欠陥補正法 DC ($k=3$)}（§\ref{sec:defect_correction_ppe}）を適用し，
        Dirichlet テストで $\Ord{h^7}$ のスーパー収束を達成する．
        単体検証 U6 では HFE 1D が $\Ord{h^{5.91}}$，
        2D 帯領域が max $11.35$ / med $8.72$ の effective slope
        （丸め床を含む診断値）を示し，
        DC 残差 $7\times10^{-6}$（$\rho_l=1000$, $\omega=0.9$, $k=3$ 最良）を実測した．
        統合検証 V6 において密度比 sweep $\rho_l/\rho_g\le833$ で
        $|\Delta p_{\mathrm{corr}}|/(\sigma/r)=1.07\times10^{-1}$ ($N=32$) を確認し，
        変密度 PPE 収束の密度比非依存性を実証した
        （pressure-jump 補正 diagnostic として測定；Laplace 絶対圧誤差ではない）．
        高密度比下の整合は pressure-jump 型 FCCD-PPE
        （§\ref{sec:split_ppe_pressure_jump_formulation}）を標準経路とする．
        ただし非一様格子では，座標変換だけで affine jump の
        $B_f(j)$ が自動的に整うわけではなく，
        $B_f^{\mathrm{nu}}=s_fj_f/H_f$ を PPE RHS と corrector が共有する
        face-jump 演算子として接続する必要がある．
\end{description}

\paragraph{第三の柱：Balanced-Force 七原則と FCCD 面ジェット sub-system}\hfill\null

\begin{description}
  \item[\textbf{第\ref{sec:collocate}章：コロケート格子・圧力--速度連成・BF 設計}]
        コロケート格子でのチェッカーボード不安定のメカニズム（$2\Delta x$ モード解析）を示し，
        変密度 IPC Projection の Helmholtz 分解導出（§\ref{sec:pressure}）と
        BF 失敗モード F-1..F-5 の分類を整理した．
        \textbf{Balanced-Force 七原則 P-1..P-7}（§\ref{sec:balanced_force}）として，
        位置共起・系統勾配共有・幾何経路統合・$\beta_f$ 統合・曲率非 CCD 則・
        ジャンプ補正・BF sub-system 分離 を体系化し，
        \textbf{FCCD-BF sub-system}（§\ref{sec:fccd_bf_sub_system}）が
        圧力勾配・表面張力・HFE 上流点・GFM ジャンプの 4 経路で
        単一の面足跡を共有することで原則群を実装する．
        単体検証 U7 で BF 静止液滴 Laplace 圧誤差 $0.61\%$ ($N=128$)，
        面 $\mu$ 補間 slopes 1.91/1.81/1.91 を確認した．
\end{description}

\paragraph{統合層：三本柱を結ぶ非一様格子・統合アルゴリズム・ブートストラップ}\hfill\null

\begin{description}
  \item[\textbf{第\ref{sec:grid}章：界面適合非一様格子と CCD/FCCD/Ridge--Eikonal 非一様拡張}]
        Level Set 値に連動したグリッド密度関数 $\omega(\phi)$ により
        界面近傍を自動的に高解像度化する非一様テンソル積格子を定式化し（§\ref{sec:grid_gen}），
        座標変換メトリック係数の CCD 評価と ALE 効果の簡略化条件を整理した．
        固定非一様格子では，セル中心面（$\theta=1/2$）と高次メトリクス構築を満たす範囲で
        FCCD は $\Ord{H^4}$ 設計を保ち（実測 $\Ord{h^{4.00/4.09}}$），
        Ridge--Eikonal は D1--D4 の 4 補正項と壁閉包 W
        （§\ref{sec:ridge_eikonal_nu_grid}：$\sigma_{\mathrm{eff}}$ 空間スケーリング・
        物理空間 Hessian・非一様 FMM・$\varepsilon$ 拡幅空間依存・contact-line seed/pinning）
        で非一様格子と壁面接触界面に拡張される．
        単体検証 U3 では GCL 残差 $2.13\times10^{-13}$（機械零）を達成した．

  \item[\textbf{第\ref{sec:algorithm}章：統合アルゴリズム（7 ステップ）}]
        NS 方程式の各項と離散ソルバの対応関係（演算子マッピング），
        \textbf{7 ステップ時間ループ}（CLS 移流→再初期化→物性更新→曲率→
        運動量予測子→PPE→IPC 補正子）を本稿全体で整理した．
        分相 PPE + HFE + DC 経路は固定格子・射影後面速度閉包には統合済みだが，
        Mode 2 再格子化は離散 GCL，ALE 保存更新，全変数リマップ，
        界面面データ再構築を追加ゲートとする．

  \item[\textbf{第\ref{sec:bootstrap}章：DCCD ブートストラップと Pure-FCCD DNS}]
        起動時の格子↔界面循環依存を解決する
        \textbf{6 段階ブートストラップ}（仮一様 SDF →非一様格子→ CLS 初期化→
        CCD 演算子→非一様 SDF 再初期化→ Young--Laplace 初期圧力）を定式化し，
        さらに研究級鋭利界面極限の構成として
        \textbf{Pure-FCCD DNS 4 段階アーキテクチャ}（§\ref{sec:pure_fccd_dns}）を提示した．
\end{description}

% ─────────────────────────────────────────────────────────────────────────────
\subsubsection{精度バランスの整理}
% ─────────────────────────────────────────────────────────────────────────────

本稿の数値手法における各コンポーネントの精度は，
単体検証 §\ref{sec:numerical_integrity}（表~\ref{tab:U_summary}）と
統合検証 §\ref{sec:verification}（表~\ref{tab:v_accuracy_summary}）で実測値が確定している．
ここでは，三本柱設計が達成した精度バランスを技術的貢献の観点から振り返る．

\paragraph{空間精度（U1, U2, U3, U6 で実測）}
CCD 内点 $\Ord{h^6}$（U1-a 実測 $\Ord{h^{6.0}}$）・大域 $L^2$ では $\Ord{h^{5/2}}$（境界 BC 律速；§\ref{sec:CCD}），
UCCD6 のスーパー収束 $\Ord{h^7}$（U1-d 実測 $\Ord{h^{7.00}}$；§\ref{sec:uccd6_def}），
FCCD 面ジェット $\Ord{h^4}$（U1-c 実測 $\Ord{h^{4.00/4.00}}$；§\ref{sec:fccd_def}），
HFE Hermite 場延長 1D $\Ord{h^{5.91}}$ / 2D effective slope max $11.35$ med $8.72$（U6-c；§\ref{sec:hfe}），
	DC k=3 補正後の分相 PPE $\Ord{h^{6.90}}$（U2，U6-b；§\ref{sec:defect_correction_ppe}）が
それぞれ単体で実測されている．

\paragraph{PPE 設計（本稿の標準 headline 設計）}
本稿の PPE 標準設計は \textbf{分相 PPE + pressure-jump 型 FCCD-PPE + 毛管 cochain 値域射影 + DC k=3} であり，
§\ref{sec:algorithm} の統合アルゴリズム 7 ステップに完全に組み込まれている．
旧来の「FD 直接法 PPE $\Ord{h^2}$ + CCD 勾配 $\Ord{h^6}$」併用は本稿の標準ではない．
密度比 sweep V6（§\ref{sec:density_ratio_sweep}）では，全 8 ケースが安定し，
$u_\infty^{\mathrm{end}}\le1.3\times10^{-9}$，
$|\Delta V_\psi|/V_{\psi,0}\le1.2\times10^{-16}$ を確認した．
CSF $\Ord{\varepsilon^2}$ 残存モデル誤差は pressure-jump 型 FCCD-PPE
（§\ref{sec:split_ppe_pressure_jump_formulation}）により部分的に解消され，
標準静止ゲートでは $N=32,T=10$ で
$K_{\max}=1.718\times10^{-6}$，$|\Delta V|/V_0\le2.919\times10^{-15}$，
$D=0$ を確認した（表~\ref{tab:V3_production_static_droplet}）．
完全 GFM 整合（Kang ら~\cite{Kang2000} の Ghost Fluid Method）は今後の課題である．

\paragraph{時間精度（U8, V7 で実測）}
TVD-RK3 $\Ord{\Delta t^{3.01}}$，AB2+IPC $\Ord{\Delta t^{2.04}}$，CN ADI $\Ord{\Delta t^{2.00}}$
（U8-a, U8-b；§\ref{sec:time_int}）の単体精度を達成している．
二相 V7（§\ref{sec:imex_bdf2_twophase_time}）の effective slope $1.59$ は
capillary pressure-jump/projection 界面帯を含む coupled-stack 律速であり，条件付き $\triangle$
（Type-D, pressure-jump capillary stack に支配され BDF2 単体 $2.0$ に届かない regime）
として記録した．
さらなる高次化（Guermond--Shen 2006 等の高次 Projection）は今後の課題である．

これらの特性は意図的な設計選択であり，
三本柱（CCD 高次群／分相 PPE+DC+CCD 勾配／BF+FCCD 面ジェット sub-system）の整合により，
寄生流れの抑制（V5 で CCD $u_\infty^{\mathrm{end}}=1.93\!\times\!10^{-2}$ @ $\rho_r=1$, $N=32$（FD 比 $1/5.44$ 補助）；§\ref{sec:spurious_multistep}）と
	高密度比下の pressure-jump PPE スタック安定性（V6 で $\rho_r \le 833$ までブローアップなし，$u_\infty^{\mathrm{end}}\le1.3\!\times\!10^{-9}$，$|\Delta V_\psi|/V_{\psi,0} \le 1.2\!\times\!10^{-16}$；§\ref{sec:density_ratio_sweep}）を同時に達成している．

% ═══════════════════════════════════════════════════════════════════════════════
